% GENERATED FILE. Edit canonical lessons, then run npm run generate:books.
% canonical-source-hash: fnv1a64:82e813d0c48cafb8
% canonical-lessons: UR-C120-halva, UR-C120-khir, UR-C120-chatni, UR-C120-kabab, UR-C120-biryani

\chapter{Halwa, Rice pudding, Chutney, Kebab, Biryani}
\label{ch:ur-halwa-rice-pudding-chutney-kebab-biryani}
\hlchaptermodality{120}

\begin{chapteropening}
\textbf{By the end of this chapter:} \emph{I can say halwa, rice pudding, a chutney, a kebab and biryani.}

Complete the last lesson of chapter 120: I can say halwa, rice pudding, a chutney, a kebab and biryani.
\end{chapteropening}

\section[\texorpdfstring{\ur{حلوہ} (halvā)}{حلوہ (halvā)}]{\ur{حلوہ} (halvā) — halwa, a sweet}

\label{lesson:UR-C120-halva}



\begin{quote}
Before the new one: say the Urdu for a paratha, then the Urdu for a samosa.
\end{quote}

\subsection*{You'll want to know: \ur{حلوہ}}
\textbf{\ur{حلوہ}} — \emph{halvā} — \textquotedblleft{}halwa, a sweet\textquotedblright{}.

\begin{cousinweb}
\textbf{\ur{حلوہ}} (\emph{halvā}) is a loan from Arabic \emph{ḥalwā}, \textquotedblleft{}a sweetmeat\textquotedblright{}, which reached Urdu through Persian. Its Arabic root is \emph{ḥ-l-w} (Lane).
\end{cousinweb}

\begin{grammarlens}[title={What you've built}]
One more word for this chapter.
\end{grammarlens}

\subsection*{Your turn}
\begin{itemize}
\raggedright
  \item \emph{Say it:} \emph{halvā}
  \item \emph{Say it:} \emph{halvā}, once more
  \item \emph{Say it:} say \emph{halvā}
  \item \emph{Recall:} say the Urdu for a paratha, then the Urdu for a samosa, then say \emph{halvā} again
\end{itemize}

\begin{tcolorbox}[breakable,colback=teal!4,colframe=teal!35!black,title={Before you move on}]
What does \ur{حلوہ} mean? (\textquotedblleft{}Halwa, a sweet\textquotedblright{}.) Say it once more.
\end{tcolorbox}
\section[\texorpdfstring{\ur{کھیر} (khīr)}{کھیر (khīr)}]{\ur{کھیر} (khīr) — rice pudding}

\label{lesson:UR-C120-khir}



\begin{quote}
Before the new one: say the Urdu for a samosa, then the Urdu for halwa.
\end{quote}

\subsection*{You'll want to know: \ur{کھیر}}
\textbf{\ur{کھیر}} — \emph{khīr} — \textquotedblleft{}rice pudding\textquotedblright{}.

\begin{cousinweb}
\textbf{\ur{کھیر}} (\emph{khīr}) is a tadbhava, a word inherited from Sanskrit and reshaped by sound change. It comes from \emph{kṣīrī}, \textquotedblleft{}a dish prepared with milk\textquotedblright{} (CDIAL 3705). It is built on \emph{kṣīra}, \textquotedblleft{}milk\textquotedblright{}.
\end{cousinweb}

\begin{grammarlens}[title={What you've built}]
One more word for this chapter.
\end{grammarlens}

\subsection*{Your turn}
\begin{itemize}
\raggedright
  \item \emph{Say it:} \emph{khīr}
  \item \emph{Say it:} \emph{khīr}, once more
  \item \emph{Say it:} say \emph{khīr}
  \item \emph{Recall:} say the Urdu for a samosa, then the Urdu for halwa, then say \emph{khīr} again
\end{itemize}

\begin{tcolorbox}[breakable,colback=teal!4,colframe=teal!35!black,title={Before you move on}]
What does \ur{کھیر} mean? (\textquotedblleft{}Rice pudding\textquotedblright{}.) Say it once more.
\end{tcolorbox}
\section[\texorpdfstring{\ur{چٹنی} (chaṭnī)}{چٹنی (chaṭnī)}]{\ur{چٹنی} (chaṭnī) — a chutney}

\label{lesson:UR-C120-chatni}



\begin{quote}
Before the new one: say the Urdu for halwa, then the Urdu for rice pudding.
\end{quote}

\subsection*{You'll want to know: \ur{چٹنی}}
\textbf{\ur{چٹنی}} — \emph{chaṭnī} — \textquotedblleft{}a chutney\textquotedblright{}.

\begin{cousinweb}
\textbf{\ur{چٹنی}} (\emph{chaṭnī}) is usually tied to \emph{chāṭnā}, \textquotedblleft{}to lick\textquotedblright{}. CDIAL lists that verb, and \emph{chaṭnā}, \textquotedblleft{}to be tasted\textquotedblright{}, made from it, under a reconstructed \emph{*caṭṭ}, \textquotedblleft{}lick, taste\textquotedblright{} (CDIAL 4573), but it does not list \emph{chaṭnī} itself. English took the word as \emph{chutney}.
\end{cousinweb}

\begin{grammarlens}[title={What you've built}]
One more word for this chapter.
\end{grammarlens}

\subsection*{Your turn}
\begin{itemize}
\raggedright
  \item \emph{Say it:} \emph{chaṭnī}
  \item \emph{Say it:} \emph{chaṭnī}, once more
  \item \emph{Say it:} say \emph{chaṭnī}
  \item \emph{Recall:} say the Urdu for halwa, then the Urdu for rice pudding, then say \emph{chaṭnī} again
\end{itemize}

\begin{tcolorbox}[breakable,colback=teal!4,colframe=teal!35!black,title={Before you move on}]
What does \ur{چٹنی} mean? (\textquotedblleft{}A chutney\textquotedblright{}.) Say it once more.
\end{tcolorbox}
\section[\texorpdfstring{\ur{کباب} (kabāb)}{کباب (kabāb)}]{\ur{کباب} (kabāb) — a kebab}

\label{lesson:UR-C120-kabab}



\begin{quote}
Before the new one: say the Urdu for rice pudding, then the Urdu for a chutney.
\end{quote}

\subsection*{You'll want to know: \ur{کباب}}
\textbf{\ur{کباب}} — \emph{kabāb} — \textquotedblleft{}a kebab\textquotedblright{}.

\begin{cousinweb}
\textbf{\ur{کباب}} (\emph{kabāb}) is a loan from Arabic \emph{kabāb}, \textquotedblleft{}meat cut up and roasted\textquotedblright{}, which reached Urdu through Persian. Lane lists it under the root \emph{k-b-b}, and records that some Arab scholars took the word to be Persian.
\end{cousinweb}

\begin{grammarlens}[title={What you've built}]
One more word for this chapter.
\end{grammarlens}

\subsection*{Your turn}
\begin{itemize}
\raggedright
  \item \emph{Say it:} \emph{kabāb}
  \item \emph{Say it:} \emph{kabāb}, once more
  \item \emph{Say it:} say \emph{kabāb}
  \item \emph{Recall:} say the Urdu for rice pudding, then the Urdu for a chutney, then say \emph{kabāb} again
\end{itemize}

\begin{tcolorbox}[breakable,colback=teal!4,colframe=teal!35!black,title={Before you move on}]
What does \ur{کباب} mean? (\textquotedblleft{}A kebab\textquotedblright{}.) Say it once more.
\end{tcolorbox}
\section[\texorpdfstring{\ur{بریانی} (biryānī)}{بریانی (biryānī)}]{\ur{بریانی} (biryānī) — biryani}

\label{lesson:UR-C120-biryani}



\begin{quote}
Before the new one: say the Urdu for a chutney, then the Urdu for a kebab.
\end{quote}

\subsection*{You'll want to know: \ur{بریانی}}
\textbf{\ur{بریانی}} — \emph{biryānī} — \textquotedblleft{}biryani\textquotedblright{}.

\begin{cousinweb}
\textbf{\ur{بریانی}} (\emph{biryānī}) is a loan, usually traced to Persian \emph{biryān}, \textquotedblleft{}fried, roasted\textquotedblright{}.
\end{cousinweb}

\begin{grammarlens}[title={What you've built}]
One more word for this chapter.
\end{grammarlens}

\subsection*{Your turn}
\begin{itemize}
\raggedright
  \item \emph{Say it:} \emph{biryānī}
  \item \emph{Say it:} \emph{biryānī}, once more
  \item \emph{Say it:} say \emph{biryānī}
  \item \emph{Recall:} say the Urdu for a chutney, then the Urdu for a kebab, then say \emph{biryānī} again
\end{itemize}

\begin{tcolorbox}[breakable,colback=teal!4,colframe=teal!35!black,title={Before you move on}]
What does \ur{بریانی} mean? (\textquotedblleft{}Biryani\textquotedblright{}.) Say it once more.
\end{tcolorbox}
